\begin{abstract}
Hierarchical summarization lets practitioners process documents that exceed an LLM's context window, but it also obscures whether the resulting summaries still encode the task value that analysts care about. We provide a probabilistic process-control framework for decomposable tasks (event extraction, counting, coding, and guided summaries with explicit rubrics) that estimates fidelity without rereading the full text.

This paper makes three contributions. \textbf{First}, we introduce Oracle-Preserving Summarization (OPS), built on a trio of consistency conditions---Sufficiency, Merge Consistency, and Idempotence---that test only the inputs accessible at each node. When these checks hold, the ``oracle'' value is preserved in expectation, and we can estimate the pipeline's empirical violation rate by sampling a fixed budget of leaves, internal merges, and leaf boundaries per audit run. The same signals serve as optimization objectives for DSPy-style prompt programs.

\textbf{Second}, we extend the framework to preference learning. We prove that when downstream supervision depends only on the oracle (e.g., DPO or GRPO loss), training on audited summaries is information-equivalent to training on the raw documents. Our bounds cover DPO, GRPO with Plackett-Luce ranking (generalizing to $k$-wise comparisons), and GRPO-RL objectives including DeepSeek-R1 style training. A unified gap bound shows: $|\mathcal{L}_{\text{orig}} - \mathcal{L}_{\text{summary}}| \le L \cdot \Delta_R$, where $L$ is a method-specific Lipschitz constant and $\Delta_R$ is expected distortion.

\textbf{Third}, we provide machine-verified proofs of all theoretical results. The formalization now comprises 50,000+ lines of Lean~4 code across 123 files with a single justified axiom (Expected Lipschitz from Random Utility Models). We release an open-source Python package, \textsc{ThinkingTrees}, implementing the full pipeline: hierarchical tree construction, probabilistic auditing with Hoeffding bounds, DSPy-first optimization (GEPA, MIPRO, Bootstrap), and preference learning integration.

We demonstrate the framework on two applications in political science: RILE scoring of party manifestos from the Manifesto Project, and policy position extraction from the Congressional Record. The result is a practical recipe for building, auditing, and tuning long-context pipelines that provide statistical guarantees about the process without requiring end-to-end verification of every document.
\end{abstract}
